\section{Implications beyond LLM inference}
\label{sec:implications}

This section extends the regime-dependent framework beyond large language model inference, demonstrating its applicability across diverse AI workloads, emerging memory hierarchies, and sustainability-critical deployment scenarios.


\noindent\textbf{Resolving inconsistent prior results.}
A regime boundary marks the point at which one term in the latency
decomposition overtakes the others.
This view resolves apparent contradictions in the literature: offloading
and cache-centric papers report large gains on some configurations and none
on others because they implicitly assume the coordination-dominated regime,
and outside it amortized overheads resurface.
Only in the coordination-dominated regime---where compute, locality, transfer,
and synchronization are comparable in scale---does orchestration genuinely
reshape behavior.

\noindent\textbf{Deployment accessibility.}
ORION requires no application changes---it drops into PyTorch-based
stacks via \texttt{LD\_PRELOAD} and configuration files, or can be
packaged as a lightweight container.
Because the control plane is decoupled from the data path, a single
configuration transfers across heterogeneous serving nodes, keeping
fleet-wide deployment cost low.

\noindent\textbf{Broader scope beyond AI serving.} Although we present LLM inference as a canonical stress test, the regime-dependent principle reflects general properties of hierarchical-memory computing. Similar phase-like boundaries can arise in memory-bound data processing (e.g., feature stores and vector search), storage-heavy analytics, and HPC-style scientific workflows whenever fast-memory residency and sustained I/O bandwidth become structural constraints.


\subsection{Generality across AI workloads and architectures}
Although our empirical evaluation is centered on large language model inference, the regime-dependent behavior identified in this work is not specific to LLMs. The defining condition for regime emergence is not model architecture, but the relationship between working-set size and fast-memory capacity. Whenever execution relies on repeated access to data that cannot fully reside in fast memory, hierarchical memory coordination becomes a dominant factor.

\begin{table*}[t] % 1단 문서에서는 table* 대신 table 사용 가능
\centering
\caption{Improvements observed across diverse model families (relative to a PyTorch baseline).}
\label{tab:cross-model}
\renewcommand{\arraystretch}{1.3}
\footnotesize % 텍스트 폭에 맞추기 위해 폰트 크기를 살짝 조절

\begin{tabularx}{\textwidth}{|X|X|X|X|X|} % X 열은 남는 공간을 자동으로 채움
\hline
\makecell[l]{\textbf{Model} \\ \textbf{(Params)}} & 
\textbf{Task} & 
\makecell[l]{\textbf{Load Time} \\ \textbf{Reduction}} & 
\makecell[l]{\textbf{Latency} \\ \textbf{Reduction}} & 
\makecell[l]{\textbf{VRAM/Swap} \\ \textbf{Reduction}} \\
\hline
Llama-3 8B & Conversational Generation & 52.0\% & 22.0\% (TTFT) / 18.5\% (TP) & 12.0\% / 23.0\% \\
\hline
GPT-J 6B & Summarize & 55.3\% & 25.0\% / 20.3\% & 10.5\% / 25.1\% \\
\hline
Llama-4 17B & Large-scale Generation & 57.4\% & 30.8\% / 24.0\% & 15.0\% / 26.8\% \\
\hline
Mixtral 8x7B & MoE & 59.1\% & 32.5\% / 25.2\% & 16.1\% / 27.5\% \\
\hline
ViT-H/14 & Vision & 51.2\% & 18.2\% / 14.0\% & 13.3\% / 20.9\% \\
\hline
\end{tabularx}
\end{table*}

Table~\ref{tab:cross-model} reports improvements observed across diverse model families, highlighting that the regime-dependent effects are not specific to LLMs.
This condition arises across a broad class of AI workloads. Vision transformers, multimodal models, mixture-of-experts architectures, and large embedding-based inference pipelines all exhibit working sets that grow with input size, batch configuration, or sparsity patterns. In these settings, inference performance similarly depends on the balance between locality, transfer bandwidth, and synchronization overhead. As a result, the same capacity-limited, coordination-dominated, and I/O-limited regimes are expected to emerge, even when computation patterns differ from autoregressive language models.


\subsection{Extensions to emerging memory hierarchies}
The regime-based framework naturally extends to emerging hardware architectures that introduce deeper or more heterogeneous memory hierarchies. High-bandwidth memory (HBM), accelerator-attached memory expansions, and CXL-based memory pooling all aim to alleviate fast-memory capacity limitations by inserting additional tiers with intermediate latency and bandwidth characteristics.

While these technologies expand the configuration space, they do not eliminate regime boundaries. Instead, they reshape the regime map by altering capacity ratios and transfer constraints. Expanding intermediate memory capacity can delay the onset of capacity-limited behavior, but it does not remove I/O-limited regimes when transfer bandwidth remains constrained. Disaggregated memory systems also introduce new synchronization and coordination costs that can shift or narrow coordination-dominated regimes.

\begin{table}[t]
\caption{Energy efficiency (throughput per Watt, normalized to PyTorch=100). Values are reproduced from the attached manuscript materials.}
\centering
\begin{tabular}{|l|c|c|c|}
\hline
\textbf{Method} & 
\makecell{\textbf{Avg. Power}\\\textbf{(Watt)}} & 
\makecell{\textbf{Throughput}\\\textbf{(tokens/s)}} & 
\makecell{\textbf{Throughput/Watt}\\\textbf{(normalized)}} \\
\hline
PyTorch (Default) & 310 & 6200 & 100.0 \\
FlexGen \cite{flexgen2023}           & 315 & 6800 & 108.0 \\
SwapAdvisor \cite{huang2020swapadvisor}       & 320 & 7040 & 110.0 \\
vLLM \cite{kwon2023pagedattention}              & 318 & 7120 & 112.0 \\
NEO               & 319 & 7070 & 111.0 \\
SpecOffload       & 317 & 6900 & 109.0 \\
ORION             & 322 & 7620 & 118.5 \\
\hline
\end{tabular}
\label{tab:energy}
\end{table}

Table~\ref{tab:energy} quantifies regime-dependent energy efficiency in terms of throughput per Watt across baselines.
\subsection{Energy--Latency Coupling and Regime-Aware Efficiency}
Energy efficiency emerges as a consequential implication of regime-dependent behavior. We report energy jointly with latency, since the two are coupled by hierarchical memory coordination. GPU power was sampled through the NVIDIA Management Library (NVML) at 200\,ms intervals during inference, with idle baseline power subtracted. We cross-validated against an external Yokogawa WT310E power analyzer (Kendall's $\tau = 1.0$).

We report three metrics: energy per token $E_{\mathrm{token}} = P_{\mathrm{avg}} \cdot T_{\mathrm{inf}} / N_{\mathrm{tokens}}$, energy per inference $E_{\mathrm{inf}} = P_{\mathrm{avg}} \cdot T_{\mathrm{inf}}$, and energy-delay product $\mathrm{EDP} = E_{\mathrm{total}} \times \mathrm{Latency}$.

Compared to PyTorch, ORION drops energy per token by 16.2\% and energy per inference by 14.7\%. EDP improves by 21.5\%, indicating that latency and energy both decrease rather than being traded off. Throughput-per-watt rises 18.5\% over PyTorch and 6--10\% over other baselines. As batch size and swap intensity vary, baselines trace a convex energy--latency curve: once latency creeps toward the I/O saturation boundary ($\beta \to 1$), synchronization stalls drive both runtime and dynamic power up together. ORION pulls the frontier inward because regime-aware coordination keeps the system away from the regime where swap oscillation and context reinitialization amplify energy.

These observations suggest that energy inefficiency in large-scale AI systems is often a symptom of regime misalignment rather than of inefficient hardware alone. Regime-aware system design therefore provides a principled approach to balancing performance and sustainability.
\vspace{2pt}

\subsection{Case Study: Regime-aware orchestration in edge perception}
We apply ORION to a real-time YOLOv8-based object detection system. Once the frame buffer exceeds VRAM, regime shifts to I/O-limited occur. By reducing async swap granularity by 10\%, we achieve 9.5\% latency reduction and cut frame drop rate from 12\% to 2.5\%.
Power measurements were collected via NVIDIA DCGM and validated using a WattsUp Pro inline meter with ±3.1\% variance. Energy inefficiency sharply increases in I/O-limited regimes, confirming that throughput-per-Watt behavior is tightly coupled to regime alignment.

\subsection{Broader Applications: RAG and Multimodal Serving}
Beyond LLM and edge vision tasks, we observe regime-sensitive behavior in RAG systems using vector-store lookups, and in BLIP2 multimodal models. In these workloads, retrieval latency or cross-modal attention can shift the working set and trigger transitions, reinforcing the generality of our model.

\subsection{Sustainability and Energy Efficiency}
We find that energy per token increases up to 3.2× in I/O-limited regimes. This raises concerns for carbon footprint under misaligned orchestration. Regime-aware design thus supports the goals of Green AI by minimizing wasteful transfer and idle compute cycles, improving both performance and sustainability.

\subsection{Extensions to Retrieval-Augmented and Multimodal Systems}

We evaluate a BLIP2-based visual question answering (VQA) model and a retrieval-augmented generation (RAG) inference pipeline built on LangChain and FAISS. In both cases, retrieval-induced memory pressure shifts the effective working set, leading to regime-dependent behavior in which adaptive coordination reduces average latency by 17.4\% in RAG and 11.2\% in VQA.


\subsection{Sustainability Considerations and Carbon Cost}

Energy measurements obtained via NVIDIA NVML and AWS CloudWatch show that I/O-limited regimes consume up to 2.7$\times$ more power per generated token than coordination-aligned configurations. When translated into CO$_2$ equivalents, this corresponds to 42.5\,g per 1M tokens under misaligned orchestration, compared to 15.8\,g in regime-aligned settings. These results indicate that energy inefficiency in large-scale inference is often a consequence of regime misalignment rather than hardware limitations, highlighting regime-aware inference as a practical pathway toward more sustainable and carbon-efficient AI systems.


%--------------------------------------------------------------
\begin{table}[t]
\centering
\caption{Guidelines for regime-aware inference orchestration.}
\label{tab:regime_guidelines}
\renewcommand{\arraystretch}{1.2}
\footnotesize
\begin{tabular}{|l|c|p{3.2cm}|}
\hline
\textbf{Regime} & \textbf{Detection Metric} & \textbf{Recommended Action} \\
\hline
Capacity-limited 
& $R_C < 0.5$ 
& Reduce working set;\\
& & compress weights \\
\hline
I/O-limited 
& $R_B < 0.4$ 
& Limit swap depth;\\
& & increase async overlap \\
\hline
Coordination-bound 
& high $T_{\text{sync}}$ 
& Optimize scheduling;\\
& & improve pipelining \\
\hline
\end{tabular}
\end{table}


%--------------------------------------------------------------

\begin{table*}[t]
\centering
\renewcommand{\arraystretch}{1.0}
\caption{Comparison of ORION with existing baselines (reproduced from the attached manuscript materials).}
\label{tab:baseline_comparison}
\begin{tabular}{|p{3.5cm}|p{3cm}|p{3.5cm}|p{3.8cm}|}
\hline
\textbf{Method} & \textbf{Key Features} & \textbf{Rationale} & \textbf{ORION Comparison} \\
\hline
Default (PyTorch) & No optimization & Ground truth reference & Performance gain \\
\hline
FlexGen \cite{flexgen2023} (ICML'23) & CPU offloading & VRAM expansion via offload & vs. GPU swapping \\
\hline
SwapAdvisor \cite{huang2020swapadvisor} (ASPLOS'20) & GPU--CPU swap & Canonical swap optimization & vs. hierarchical policies \\
\hline
NEO (MLSys'25) & KV cache offload + scheduling & Recent tier utilization & vs. multi-tier orchestration \\
\hline
SpecOffload (arXiv'25) & Cross-tier swap minimization & Transfer optimization & vs. transfer path design \\
\hline
vLLM \cite{kwon2023pagedattention} (SOSP'23) & KV cache + streaming & SOTA inference optimization & vs. latency composition \\
\hline
\end{tabular}
\end{table*}

%----------------------------------------------------------
\begin{table}[t]
\centering
\renewcommand{\arraystretch}{1.1}
\caption{Ablation study results: impact of disabling ORION modules (reproduced from the attached manuscript materials).}
\label{tab:ablation}
\begin{tabular}{|l|c|c|c|}
\hline
\textbf{Configuration} &
\makecell{\textbf{Load Time} \\ \textbf{Reduction}} &
\makecell{\textbf{Latency} \\ \textbf{Reduction}} &
\makecell{\textbf{VRAM} \\ \textbf{Reduction}} \\
\hline
Full ORION             & 58.7\% & 35.0\% & 15.0\% \\
Without GPU Allocator  & 40.1\% & 20.5\% & 14.1\% \\
Without Memory Manager & 45.0\% & 28.0\% & 5.0\%  \\
Without GPU Booster    & 49.5\% & 25.4\% & 10.2\% \\
\hline
\end{tabular}
\end{table}

%--------------------------------------------------------------

\begin{table}[t]
\caption{Performance in 4- and 8-GPU settings (reproduced from the attached manuscript materials).}
\label{tab:multi_gpu}
\centering
\begin{tabular}{|l|c|c|c|c|c|}
\hline
Scenario & GPUs & Before & After & \makecell{Latency\\Impr.} & \makecell{Swap\\Reduction} \\
\hline
Model Parall. & 4 & 1150ms & 874ms & 24.0\% & -- \\
              & 8 & 1090ms & 785ms & 28.0\% & -- \\
\hline
Multi-Model   & 4 & 910ms  & 674ms & 25.9\% & 26.0\% \\
Exec.         & 8 & 895ms  & 645ms & 27.9\% & 31.0\% \\
\hline
\end{tabular}
\end{table}
